\chapter{Introduction}
\label{ch:intro}

Industrial manipulators are built to follow a commanded trajectory precisely,
and inside a structured cell that is sufficient. In surface-finishing
applications, however, the actual surface pose may differ from the programmed
pose because of placement tolerances: the surface is positioned by hand, the
tool is held in a gripper, and the residual angle between the two is not known
before contact. Under rigid position control, that initial angular deviation
cannot be absorbed by the motion and therefore generates high contact forces.

Collaborative robots make another option available, and they account for a
growing share of industrial deployments
\citep{IFR2023WorldRobotics,Villani2018HRC}. Joint torque sensing and a
torque-level command interface let the controller specify how the robot yields
under contact, not only where it goes. Three properties of the present
experimental arrangement shape what such a controller must achieve. The grinding
face is held in a gripper and is therefore not perfectly rigid with respect to
the end-effector. The surface plane is known only through a calibration that
carries its own uncertainty. The objective is consequently not force reduction
alone, but passive angular alignment of the tool during contact, measured
against a separately calibrated physical surface.

Impedance control specifies that yielding as a dynamic relation between motion
and interaction force \citep{Hogan1985}. Cartesian impedance control imposes
the relation at the end-effector, in task space, so compliance can be shaped
per direction rather than per joint. A contact task requires this compliance to
be direction dependent. A tool moved across a surface must hold its distance to
that surface, yet a small lateral misalignment should not build a large force.
Orientation behaves the same way: when a tilted tool edge meets the surface, the
contact wrench carries a moment as well as a force, and a stiffness matrix
chosen for the contact geometry admits a small rotation instead of resisting it.
The rotation that actually occurs during contact is the quantity this thesis
measures, which makes the stiffness and damping matrices a choice about contact
geometry and not only about free-space tracking.

That choice rests on three bodies of work. The first is Cartesian impedance
control itself. The controller rests on standard serial-manipulator kinematics
\citep{Siciliano2009,Craig2005,Murray1994}: homogeneous transformations for
pose, and the Jacobian for the velocity map. Khatib established the
operational-space formulation, in which a Cartesian wrench is mapped to joint
torques through \(J^\top\) \citep{Khatib1987}; that mapping is the one adopted
here. Hogan defined the impedance relation itself, between end-effector motion
and interaction force \citep{Hogan1985}. Ott et al.\ extended Cartesian
impedance control to flexible-joint robots, where joint elasticity limits the
stiffness the controller can actually realise \citep{Ott2008}, and
Albu-Sch\"affer et al.\ built the same law on joint-torque sensing
\citep{Albu2003}. The torque-sensing interface is the one this thesis uses; the
flexible-joint extension is not adopted, since joint elasticity is not modelled
here.

The second is the practice of running such a law in real time. Previous
Cartesian impedance implementations on the Franka Emika Panda discuss
requirements such as real-time model access, null-space control, signal
treatment, and command limiting \citep{Mayr2024CartesianImpedance}. The present
controller uses the libfranka model interface and includes a selectable
null-space torque. No separate application-level saturation of the Cartesian
wrench, joint torque, or torque rate was implemented, so the configured
robot-side collision monitoring remained the only independent mechanism for
terminating the motion when its thresholds were exceeded.

The third concerns where the compliance is centred. Fasse and Broenink
formulate spatial impedance about a selected point, so that translation and
rotation couple through its placement \citep{Fasse1997SpatialImpedance}. That
construction is the theoretical basis of the compliance-centre shift used here.
The same geometry appears mechanically in remote-centre-compliance devices for
assembly \citep{Nevins1989ComputerControlledAssembly,Watson1978RCC} and, more
recently, in flexure-based and variable-compliance-centre strategies for
peg-in-hole insertion, where the compliance is arranged so that contact forces
reduce rather than amplify misalignment
\citep{Wang2019VariableComplianceCenter,Hartisch2022FlexureRCC}. Those studies
mainly consider insertion, and they fix the compliance centre either through
mechanical design or through a strategy tuned for a hole. The effect of
independently varying the virtual compliance-centre position and the
axis-specific stiffness during planar tool alignment has received less
attention. This thesis investigates that effect using a software-defined
compliance centre on a torque-controlled robot, and measures what its placement
does to tool alignment against a surface.

The problem it addresses can be stated directly. The angular relation between
the tool and the surface plane is uncertain before contact, because both the
surface placement and the tool mounting carry tolerances, and under rigid
position control that residual angle produces an undesirable contact moment
rather than a corrective motion. A tool tilted with respect to the surface
cannot be brought flat by translational compliance alone: the wrench that closes
the gap must also carry a moment. Translational compliance therefore does not by
itself define the corrective rotation. Placing the centre of compliance away
from the TCP couples translation and rotation, so that pressing the tool into
the plane also rotates it toward the plane. The placement is a
task-to-impedance mapping, not a claim that the tool pivots about that point
physically, and the response it produces is therefore read from the motion
measured during contact.

The alignment response cannot be predicted from the gain values alone, because
it also depends on the contact geometry, the lever direction, the robot
configuration, and the compliance of the tool mounting. The experiments
therefore compare the effects of commanded tilt, directional stiffness, and
compliance-centre placement under otherwise matched conditions.
\Cref{ch:theory,ch:implementation} set out the frame conventions those
comparisons rely on.

\label{sec:objectives}
Within that scope, this work investigates how the commanded tilt, the rotational
and translational stiffness of each surface tangent, and the spatial position of
a virtual centre of compliance influence contact-induced alignment. The
controller is a phase-scheduled Cartesian impedance law, written in C++ with
libfranka and Eigen and running in the \(1\,\mathrm{kHz}\) torque loop of the
Franka Emika Panda. Torques are mapped with \(J^\top\), and Coriolis terms are
compensated from the model. Each phase carries its own Cartesian gains; damping
follows from a Cartesian-inertia estimate rather than being set by hand. A
surface frame carries the normal and the two tangents \(t_1\) and \(t_2\)
explicitly.

One run proceeds through fixed phases: the tool axis is aligned, the plane is
approached to a geometric clearance, the leading feature of the rectangular
tool face is selected, that feature is ramped into the surface, and tangential
grinding motion follows. During set-up, a compliance-centre \(6\times6\)
stiffness and damping matrix is shifted to the tool centre point (\abbr{TCP})
through an adjoint congruence. Its off-diagonal blocks are the
translation--rotation coupling this thesis measures.

The principal contributions are of two kinds. The following were implemented and
verified against the running controller:

\begin{itemize}
  \item A frame-consistent real-time Cartesian impedance controller for the
        7-\abbr{DOF} Panda.
  \item A spatial compliance-centre formulation within that controller, in which
        the \(6\times6\) stiffness and damping are shifted to the TCP.
  \item A calibrated contact methodology, in which the surface plane and the
        grinding-face axis are measured independently of each other.
\end{itemize}

The following were established experimentally:

\begin{itemize}
  \item Axis-specific measurement of the rotational and translational stiffness
        contributions to contact-induced alignment.
  \item Identification of the compliance-centre position as the parameter with
        the largest measured influence on alignment within the tested range,
        including the direction dependence of that influence: the lever
        direction that improves alignment is opposite for the two surface
        tangents.
\end{itemize}

One part of the implementation is deliberately excluded from both lists. The
projected joint damping and the two-point smallest-singular-value bias for the
redundant seventh joint complete the controller architecture, but no matched
comparison isolates their effect, so they are not claimed as experimentally
established. Full operational-space inverse dynamics, adaptive impedance,
online surface estimation, and claims concerning human safety are outside the
scope.

The remainder of the thesis follows the order of that investigation.
Chapter~\ref{ch:theory} sets out the impedance law, the frame conventions, and
the compliance-centre shift; Chapter~\ref{ch:implementation} gives their
real-time realisation on the Panda. Chapters~\ref{ch:methodology}
and~\ref{ch:results_discussion} describe the experimental methodology and
present the corresponding results. Chapter~\ref{ch:conclusion} states the
conclusions and the work that remains.
